\chapter{繰り込み、自然さ、そして階層性問題}

\lectureoverview{この講義の中心的なアイデアは、標準模型における繰り込みに関する謎から生まれた。したがって、まず繰り込みの物理的意味から始める。すなわち、当面の問いに対して小さすぎ、速すぎる自由度を消去し、その効果を有効な記述に繰り込むのである。分子がその原型を与える。次に次元解析が同じ推論をスカラー量子場の理論へと運び、そこではスカラー質量の二次的な感受性、ヒッグス場の階層性問題、フェルミオンとゲージボソンが享受する保護、そしてさらに深刻な真空エネルギーの問題が明らかになる。}

\section{問いのスケールにおける物理}

今学期に展開されるアイデアの大半は、標準模型における繰り込み---時に謎であり、時に有用な観察であった---に関する問いから生まれた。繰り込みは二つの思考習慣を組み合わせる。第一は、問われている事柄にとって無関係なほど短い距離の構造を記述から取り除くことである。第二は次元解析であり、これは難しい積分を評価する前に短距離の答えをしばしば決定する。

このような細部の消去は物理学の至る所で起こる。原子核を理解したいなら、クォークと量子色力学から始めてもよい。陽子と中性子の質量とスピン、およびそれらの有効な核力を計算した後なら、多くの目的のためにそれらのクォーク部分構造を忘れることができる。核の運動は十分に遅いので、陽子と中性子の非相対論的量子力学の方がしばしば有用な記述となる。

次の段階では、原子物理学は通常、原子核の内部構成を必要としない。核物理学はどの原子核と同位体が存在するかを教え、質量や電荷といった性質を与える。そこで我々は新たな有効な基本対象、すなわち原子核と電子から再び始める。クォーク、陽子、中性子が存在しなくなったわけではない。それらの効果は原子核が担うパラメータに圧縮されたのである。

どれだけ保持するかは次の問いに依存する。原子の計算では、換算質量に入るため原子核の質量が必要になることがあり、電荷はクーロン場を決定するため必ず必要である。通常、どのクォークが核の運動量のどの割合を担うかを知る必要はない。元素の電荷はおよそ九十種類あり、同位体を数えればさらに多くの核種が存在する。その目録は入力として取ることができる。一方、電子はまだ消去できない。その運動こそが原子構造を作り出しているからである。

原子の問題を解いた後は、同様に原子を分子の構成要素と見なしてよい。サスキンドの気まぐれな、クッキー原子、コーヒー原子、チョコレート原子、ケチャップ原子というリストは、あらゆる段階で新たな記述が始まるという真剣な点を突いている。分子の性質と力が分かれば、各段階で電子の問題を再び開くことなく、分子をより大きなティンカートイのような集合体に組み上げることができる。

それぞれの移行は粗視化である。新しい記述は絶対的な意味ではより正確でないかもしれないが、意図されたスケールでははるかに有用である。量子場の理論では、場があらゆる波長のモードを持ち、任意に短いものも含むため、この問題は避けられない。したがって、分解できないほど短い波長の効果を足し上げ、それらを新たな有効パラメータで置き換える。この操作が次元解析とともに繰り込みの物理的内容である。

\section{原型としての分子}

二つの原子を考え、それらの相対運動の簡単な記述を求めよう。各原子は重い原子核と電子の雲を含む。陽子でさえ電子よりおよそ二千倍重く、より重い原子核はスケールの分離をさらに大きくする。有用な描像は、高速で動くハエに囲まれた一対のボウリング球である。原子核はゆっくり応答し、電子状態は素早く再配置される。

原子核の位置と運動量を\(R_a\)と\(P_a\)、共通の質量を\(M\)とし、電子の座標と運動量を\(r_i\)と\(q_i\)、質量を\(m_e\)とする。\(4\pi\epsilon_0\)のような慣習的な因子を省略すると、ハミルトニアンは模式的に次の形をとる。
\begin{align}
H={}&
\frac{P_1^2}{2M}+\frac{P_2^2}{2M}
+\frac{e^2}{R_{12}}
+\sum_i\frac{q_i^2}{2m_e}
+\frac12\sum_{i\ne j}\frac{e^2}{r_{ij}}
\nonumber\\
&-\sum_i\left(
\frac{e^2}{|R_1-r_i|}
+\frac{e^2}{|R_2-r_i|}
\right).
\label{eq:molecular-hamiltonian}
\end{align}
第一行には原子核の運動エネルギー、それらの相互のクーロン反発、および電子の運動エネルギーが含まれる。第二行には電子-電子反発と電子-原子核引力が含まれる。一般の電荷を持つ原子核に対しては、対応する\(Z_a\)の因子が復元される。時間スケールの分離が問題であるため、この講義では電荷を模式的に保つ。

特徴的な運動が二つある。原子核は重くて遅い。電子は軽くて速く、ぼやけた雲のように見える。第一近似では、原子核の位置を固定し、電子に依存するすべての項を電子ハミルトニアンにまとめる。
\begin{equation}
H_{\mathrm{el}}(R_1,R_2).
\end{equation}
各固定された原子核配置に対して、
\begin{equation}
H_{\mathrm{el}}(R_1,R_2)
\psi_0(r_1,\ldots,r_N;R_1,R_2)
=
E_{\mathrm{el}}(R_1,R_2)\psi_0.
\label{eq:electronic-ground-state}
\end{equation}
を解く。最も低い電子固有値を保持する。電子は無視されたのではない。その速い運動のすべては関数
\(E_{\mathrm{el}}(R_1,R_2)\)
にまとめられている。

論理は断熱的である。一回の速い電子の調整の間に、原子核はほとんど動かないので、その座標はパラメータとして扱える。電子固有値問題を解いた後、同じ座標が再び動くことを許される。ゆっくり変化する各配置で、電子雲はその瞬間の基底状態に従う。より正確な扱いでは励起電子面とそれらの間の遷移を保持できるが、主要な面ですでに繰り込みの考え方が現れている。

\begin{figure}[htbp]
\centering
\includegraphics[width=0.72\textwidth]{lecture_01_figure_02.png}
\caption{分子ハミルトニアンは、電子基底状態問題が導入される間も黒板に残っている。式~\eqref{eq:molecular-hamiltonian}の活字の形は、このフレームで隠された項を記録している。}
\lectureframe{00:16:19}
\end{figure}

\begin{classroomqa}
\audiencequestion{原子核の運動量は単にゼロに設定されているのか。\lecturetimestamp{00:15:39}}
\lecturerresponse{一時的にすぎない。主要な近似では、原子核は電子よりはるかにゆっくり動くので、釘付けにされているように扱われる。その運動量は有効な原子核問題を解くときに復元されるので、この近似は不確定性原理に反するものではない。}
\end{classroomqa}

電子基底状態エネルギーは今や、原子核が感じるポテンシャルへの寄与である。
\begin{equation}
V_{\mathrm{eff}}(R_1,R_2)
=\frac{e^2}{R_{12}}+E_{\mathrm{el}}(R_1,R_2).
\label{eq:molecular-effective-potential}
\end{equation}
回転および並進不変性により、二つの孤立した原子に対して、それはそれらの間隔\(R_{12}\)の関数に帰着する。非常に短い距離では原子核のクーロン反発が支配する。非常に大きな距離では相互作用はほとんどない。これらの極限の間で、電子雲は引力領域と極小を作り出しうる。原子核波動関数はその極小の近くに位置し、束縛分子を記述する。

操作の順序に注意せよ。我々は経験的な引力を推測して電子を捨てたのではない。まず固定された原子核の背景の中で電子の量子問題を解いた。その後になって初めて、それらの明示的な座標を消去し、結果として得られた固有値を保持した。短距離の物理は\(V_{\mathrm{eff}}\)の形、深さ、曲率の中に生き残っている。その極小の位置が平衡間隔を定め、その曲率が小さな振動を支配し、その漸近形が束縛状態が存在するかどうかを支配する。一つの有効関数が、その電子情報のすべてをより遅い問題へと運ぶ。

こうして複雑な多電子ハミルトニアンは、新しいポテンシャルを持つ二体の原子核ハミルトニアンになった。我々は速い自由度を消去したが、その完全な低エネルギー効果は保持した。ポテンシャルは繰り込まれたのである。このよく知られたボルン--オッペンハイマーの推論は、通常その名前で導入されることはないが、まさに我々が必要とする考え方を含んでいる。
\begin{equation}
\begin{gathered}
\text{fast microscopic dynamics}\\
\Downarrow\\
\text{effective parameters for slow dynamics}
\end{gathered}
\end{equation}

\section{自然単位と場の次元}

量子場の理論では、消去されるモードは、短距離、短波長、高振動数、または高エネルギーのモードとして等価に記述されうる。図を描く前に、次元の帳尻を定めておこう。

通常の力学は長さ、時間、質量の独立な単位を用いる。素粒子物理学では
\begin{equation}
\hbar=c=1.
\end{equation}
と設定する。するとエネルギーと質量は\(E=mc^2\)により等価であり、長さと時間は\(c=1\)により等価であり、\(\hbar=1\)により運動量は長さの逆数となる。残る次元は一つだけである。
\begin{equation}
[M]=[E]=[p]=L^{-1},
\qquad
[x]=[t]=L.
\label{eq:natural-units}
\end{equation}

これは次元を消し去るのではない。時間を長さとして、質量を長さの逆数として表せるように変換因子を選ぶのである。一つの単位を指定すれば、他のすべてがそれに従う。これが、カットオフを短い距離\(\delta\)としても、高いエネルギー\(\Lambda=1/\delta\)としても記述できる理由である。それらは二つの言語で同じスケールである。

一つの実スカラー場\(\phi\)を取り、模式的に規格化したラグランジアン密度
\begin{equation}
\mathcal L=(\partial_\mu\phi)(\partial^\mu\phi)-V(\phi),
\qquad
S=\int d^4x\,\mathcal L.
\label{eq:scalar-action}
\end{equation}
を考える。

\begin{figure}[htbp]
\centering
\includegraphics[width=0.78\textwidth]{lecture_01_figure_03.png}
\caption{局所スカラーラグランジアン密度と、その時空積分としての作用。}
\lectureframe{00:26:08}
\end{figure}

ラグランジアン密度は局所的である。すなわち一つの時空点を記述する。古典場の方程式は作用を停留させることから従う。量子論では同じ作用が経路積分の位相\(\exp(iS)\)に現れる。どちらの見方でも\(S\)の次元が決定的である。\(\hbar\)を復元すれば、位相は
\(\exp(iS/\hbar)\)
であり、\(\hbar=1\)と設定した後では、\(S\)自身が無次元でなければならない。

\begin{classroomqa}
\audiencequestion{作用は\(\hbar\)の単位をもつか。\lecturetimestamp{00:26:13}}
\lecturerresponse{はい。\(\hbar\)が1に等しく設定されているので、作用はこれらの単位では無次元である。4次元の測度は次元\(L^4\)をもつので、ラグランジアン密度は次元\(L^{-4}\)をもたなければならない。}
\end{classroomqa}

運動項は今や場の次元を固定する：
\begin{align}
[S]&=1, & [\mathcal L]&=L^{-4},\\
[\partial_\mu]&=L^{-1}, &
[(\partial\phi)^2]&=L^{-2}[\phi]^2=L^{-4},\nonumber\\
&& [\phi]&=L^{-1}.
\label{eq:scalar-dimension}
\end{align}
この小さな計算が、以下に続くすべての図の紫外的な大きさを決定する。

ポテンシャルを次のように展開する。
\begin{equation}
V(\phi)=\frac{m^2}{2}\phi^2+g\phi^3+\lambda\phi^4+\cdots.
\label{eq:scalar-potential}
\end{equation}
各項は次元\(L^{-4}\)を持たなければならない。
\([\phi^n]=L^{-n}\)
であるから、
\begin{equation}
[m^2]=L^{-2},
\qquad
[g]=L^{-1},
\qquad
[\lambda]=1.
\label{eq:coupling-dimensions}
\end{equation}
四次結合定数は無次元である。このような結合定数は、そのスケール依存性を単純な長さの冪から読み取ることができないため、繰り込みにおいて特に重要である。

\begin{classroomqa}
\audiencequestion{\(g\)は質量の単位をもつか。\lecturetimestamp{00:29:52}}
\lecturerresponse{はい。因子\(\phi^3\)は次元\(L^{-3}\)をもち、一方ポテンシャル密度は次元\(L^{-4}\)をもたなければならない。したがって\([g]=L^{-1}\)であり、これは自然単位における質量である。}
\end{classroomqa}

\section{頂点と短距離伝播関数}

ファインマン則は頂点と伝播関数から構成される。項\(g\phi^3\)は\(g\)に比例する因子を持つ三本線の頂点を与え、\(\lambda\phi^4\)は\(\lambda\)に比例する因子を持つ四本線の頂点を与える。二次の項は局所的な二点挿入を与える。これは粒子が入射し、吸収され、再放出されるものとして描くことができ、その係数は\(m^2/2\)に比例する。数値的な規格化、符号、\(i\)の因子、対称因子は意図的に省略している。というのも、この講義は断面積を計算するのではなく、スケール依存性を取り出そうとしているからである。

これらの図には操作的な読み方がある。三点頂点は二個の量子を吸収して一個を放出することができ、あるいはその三本の線の交差に関連する任意の割り当てを表すことができる。四点頂点は二個の粒子が入射し二個が出ていくことを記述できる。二次のクロスは、通常は自由伝播関数に取り込まれるものの、依然として正当な二本線の頂点である。これを表示するのは、未解決のすべての二点過程がまもなくこれとの比較によって認識されるからである。

伝播関数は2つの時空点を結ぶ。その線は、量子粒子が古典的な直線経路をたどったという主張ではない。それは、\(x\)で場の励起を生成し、\(y\)でそれを除去または検出する振幅を表す：
\begin{equation}
\Delta(x-y)=\langle0|\phi(y)\phi(x)|0\rangle.
\label{eq:position-propagator}
\end{equation}
この式を右から読む。\(\phi(x)\)を真空に作用させると\(x\)に1粒子成分が生成される。\(\phi(y)\)をブラに作用させると\(y\)の成分が除去される。それらの真空行列要素は、第1の点での生成と第2の点での検出の間の重なり振幅である。記号\(x\)と\(y\)は、単に2つのデカルト座標ではなく、4元ベクトル全体を表す。

状態は物理的な次元を持たず、各場は\(L^{-1}\)を寄与する。したがって
\begin{equation}
[\Delta]=L^{-2}.
\end{equation}
質量のない場には固有の長さスケールが存在しないため、次元解析により短距離での振る舞い
\begin{equation}
\Delta(x-y)\sim\frac{1}{(x-y)^2}.
\label{eq:massless-propagator}
\end{equation}
が強制される。分母はローレンツ不変な分離の二乗であり、その特異点において通常の処方が必要である。\editorialnote{パワーカウントには、大きさが\(1/|x-y|^2\)に曖昧さなく比例するユークリッド的分離へウィック回転するのが便利である。これは以下で用いる紫外の冪を変えない。}

\(y\)が\(x\)に近づくにつれて、伝播関数は発散する。消えゆく分離を越えて伝播するこの特異性が、より複雑な図における紫外発散が組み立てられる源である。

\section{スカラー質量はいかにして繰り込まれるか}

スカラー場では、相対論的分散関係がすでに示唆するように、力学的パラメータは\(m^2\)である：
\begin{equation}
E=\sqrt{\mathbf p^2+m^2}.
\end{equation}
\(\delta\)を、我々の有効記述が分解する最短距離とする。頂点が\(\delta\)より短い距離で隔てられている過程は区別しない。1個の入りと1個の出のスカラーをもつそのような未解決の過程は、質量挿入と同じ観測可能な2点の形をもつので、有効質量に寄与しなければならない。

加速器を顕微鏡と考えよ。その最良の分解能が\(\delta\)であるなら、ある点での吸収の後に\(\delta\)よりずっと小さい距離で起こる放出は、真に局所的な事象と区別できない。我々は意図的にそのスケールにわたって頂点をぼかす。繰り込みとは、そのような分解能以下のすべての代替を対応する局所演算子の係数へ移す規則である。

\subsection{最初の四次補正}

まず単一の\(\phi^4\)頂点から始める。二本の脚は外部粒子であり、残りの二本はほぼ同じ点で放出・再吸収される短いループに合流する。一致した伝播関数をカットオフ距離にわたってならすと
\begin{equation}
\Delta(0)\ \longrightarrow\ \Delta(\delta)\sim\frac{1}{\delta^2},
\end{equation}
が得られ、したがって
\begin{equation}
\delta m^2_{(1)}\sim\frac{\lambda}{\delta^2}.
\label{eq:one-loop-scalar-mass}
\end{equation}
となる。\(m_0^2\)が入力パラメータすなわち裸のパラメータを表すなら、低分解能の二点係数は模式的に
\begin{equation}
\frac{m_{\mathrm{eff}}^2}{2}
\sim
\frac{m_0^2}{2}
+c_1\frac{\lambda}{\delta^2},
\label{eq:first-renormalized-mass}
\end{equation}
の形をとる。ここで\(c_1\)は無視した数値因子および規約依存因子を含む。\(\delta\)より短い距離を分解しない実験で測定される実験室質量はこの繰り込み量であり、\(m_0\)そのものではない。

したがって、観測される二点事象が起こる方法は二通りあり、区別できない。元の二次頂点が粒子を吸収・再放出してもよいし、四次相互作用が一時的に仮想スカラーを放出・再吸収しつつ同様にしてもよい。量子力学では振幅が加算される。ループが分解能スケールより下に入れば、その個別の履歴は有効記述から消えるが、測定される係数への寄与は消えない。

\subsection{次の四次補正}

次に二つの四次頂点を結び、一つの外部粒子が\(x\)で入り、一つが\(x+\Delta\)から出て、三本の内部伝播関数が頂点を結ぶようにする。\(\lambda\)の二乗が現れ、伝播関数は\(\Delta^{-6}\)を与える。量子振幅は、出ていく粒子が現れうる未分解のあらゆる場所について和をとらなければならないので、相対四元ベクトルを積分する：
\begin{equation}
\delta m^2_{(2)}
\sim
\lambda^2
\int_{|\Delta|\gtrsim\delta}
d^4\Delta\,\frac{1}{|\Delta|^6}.
\label{eq:two-loop-power-count}
\end{equation}

\begin{classroomqa}
\audiencequestion{頂点にある小さな印は、そこに粒子が三つしかないという意味ですか。\lecturetimestamp{00:46:35}}
\lecturerresponse{いいえ。我々は依然として四次相互作用を用いている。各頂点には四本の線が集まる。この二点図では一本の線が外部で、三本が二つの頂点の間を走る。第二の頂点にもやはり四本の線が入射する。}
\end{classroomqa}

測度は長さの四乗を与え、三本の伝播関数はマイナス六乗を与える。紫外カットオフがこの見積もりにおける唯一のスケールであるから、
\begin{equation}
\int_{|\Delta|\gtrsim\delta}
\frac{d^4\Delta}{|\Delta|^6}
\sim\frac{1}{\delta^2},
\qquad
\delta m^2_{(2)}
\sim\frac{\lambda^2}{\delta^2}.
\label{eq:two-loop-scalar-mass}
\end{equation}
ここで\(\Delta^\mu\)は四成分の積分変数であり、その時間成分と三つの空間成分について積分する。小文字の\(\delta\)は固定されたカットオフである。講義では当初、両者に同じ呼び名を用い、後に記法を訂正する。

同じ結果は四次元測度を動径的に書くことでも明らかである。角度に関する定数を除いて、
\begin{equation}
\int_{\delta}\frac{d^4\Delta}{|\Delta|^6}
\sim
\int_{\delta}\frac{R^3\,dR}{R^6}
=\int_{\delta}\frac{dR}{R^3}
\sim\frac{1}{\delta^2}.
\end{equation}
これが、紫外見積もりにおいて大きな距離の上限が重要でない理由である。積分はその下端で支配される。

より高次の図も同じパターンを繰り返す。四次頂点が多い図は\(\lambda\)のより高い冪を加えるが、次元解析は依然として二次発散する質量寄与を許す：
\begin{equation}
\begin{aligned}
m_{\mathrm{eff}}^2
={}&m_0^2
+\frac{1}{\delta^2}
\left(c_1\lambda+c_2\lambda^2+c_3\lambda^3+\cdots\right)\\
&+\text{weaker terms}.
\end{aligned}
\label{eq:scalar-mass-series}
\end{equation}

\subsection{三次相互作用と対数}

二つの三次頂点は近接した点の間に二本の伝播関数を含む。その寄与は
\begin{equation}
\delta m_g^2
\sim
g^2\int_{|\Delta|\gtrsim\delta}
\frac{d^4\Delta}{|\Delta|^4}.
\label{eq:cubic-loop}
\end{equation}
のようにスケールする。今や測度と分母は等しい冪を持つ。残る依存性は冪ではなく対数である：
\begin{equation}
\delta m_g^2
\sim
c_g g^2\log\!\left(\frac{1}{\mu\delta}\right).
\label{eq:cubic-log}
\end{equation}
基準スケール\(\mu\)が対数を無次元にする。これが講義の模式的表現\(g^2\log\delta\)の完全に次元をもった版である。対数はどんな逆冪よりもはるかにゆっくり変化するので、カットオフが短くなると\(1/\delta^2\)が支配する。

ここでも動径形が区別を明瞭にする：
\begin{equation}
\int_{\delta}^{R_{\mathrm{IR}}}
\frac{d^4\Delta}{|\Delta|^4}
\sim
\int_{\delta}^{R_{\mathrm{IR}}}\frac{R^3\,dR}{R^4}
=\log\!\left(\frac{R_{\mathrm{IR}}}{\delta}\right).
\end{equation}
測度と分母の冪が等しいと対数が生じる。分母の冪がより大きいとカットオフの逆冪が生じる。講義は角度因子を評価せよと求めているのではなく、紫外の答えを一目で見分けられるように教えているのである。

\begin{classroomqa}
\audiencequestion{対数は何を底としてとるのですか。\lecturetimestamp{00:57:42}}
\lecturerresponse{自然対数を用いる。底を変えても数値定数が掛かるだけで、その定数は\(2\)、\(\pi\)、対称因子をここでは計算していない係数に吸収される。}
\end{classroomqa}

\begin{classroomqa}
\audiencequestion{長さの対数をどうやってとるのですか、また質量補正の単位はどこから来るのですか。\lecturetimestamp{00:58:32}}
\lecturerresponse{対数は常に二つのスケールを比較する。正しくは\(1/(\mu\delta)\)のような無次元比の対数である。結合定数\(g\)は質量の一乗を持つので、\(g^2\)はすでに補正に必要な質量二乗の次元をちょうど持っている。}
\end{classroomqa}

\begin{classroomqa}
\audiencequestion{ファインマン図は振幅を与える。その絶対値の二乗は確率なのか。\lecturetimestamp{00:59:41}}
\lecturerresponse{はい、振幅が必要な状態因子と相空間因子を伴う物理的な遷移計算に置かれた場合には確率である。ここでのより直接的な役割は異なる。未解決の2点振幅は質量項と同じ局所的な形を持ち、したがって有効質量パラメータをずらすのである。}
\end{classroomqa}

\section{すべてのパラメータが有効になる}

質量図の無限族は、\(\lambda\) が小さいというだけでは消えない。摂動論はそれらを結合定数の冪で順序づけるので、逐次の項は小さくなりうるが、それでも依然として無限に多く存在する。同じ短距離の消去は、相互作用そのものも繰り込み直す。

\begin{classroomqa}
\audiencequestion{質量の繰り込みの場合でも、図は無限に多くあるのではないか。\lecturetimestamp{01:00:35}}
\lecturerresponse{ある。\(\lambda\) が十分に小さければ、より高次の冪が逐次抑えられるため、摂動級数は有用でありうる。それでも、繰り込まれたパラメータは最初の図だけでなく級数全体によって定義される。}
\end{classroomqa}

操作上、局所的な \(\lambda\phi^4\) 相互作用とは、2粒子入射・2粒子出射の振幅である。より複雑な4点過程が完全に未解決の領域の内部で起こるならば、低い分解能では元の局所的な頂点と区別できない。したがってそのような図は有効な \(\lambda\) にまとめられる。それらの主要なスケール依存性は、4次元 \(\phi^4\) 理論では典型的に対数的である。この講義ではそのベータ関数を計算しない。この例を用いて一般的な主張を行うのである。
\begin{equation}
\begin{gathered}
\text{parameters measured at low energy}\\
\ne\\
\text{parameters entered at the microscopic cutoff}
\end{gathered}
\end{equation}

四点の例が重要であるのは、ありうる誤解を取り除くからである。繰り込みは質量の発散だけに適用される特別な修復ではない。ラグランジアンで許されるすべての演算子は、同じ外部脚と同じ対称性を持つ未解決の図によって再現されうる。二点過程は二次の係数を繰り込み、四点過程は四次の係数を繰り込み、真空から真空への過程は定数項を繰り込むことになる。低エネルギーラグランジアンは、それらすべての短距離過程が記録される簡潔な台帳なのである。

\begin{classroomqa}
\audiencequestion{\(\Delta\) は長さの単位を持つので、その対数は次元的に奇妙ではないのか。また、黒板では大文字と小文字のデルタが入れ替わっていたのではないか。\lecturetimestamp{01:01:34}}
\lecturerresponse{対数は実際には別のスケールとの比を意味するので、無次元である。はい、記法は一時的に入れ替わっていた。\(\Delta^\mu\) が積分変数であり、\(\delta\) が答えに残る下限の距離カットオフである。}
\end{classroomqa}

これは分子の場合の場の理論における対応物である。われわれは見えない速い電子について解き、その結果をポテンシャルに入れた。ここでは見えない短距離過程を足し上げ、その結果を質量、結合定数、その他すべての許される局所的な係数に入れる。すべてがずれるのである。

\begin{classroomqa}
\audiencequestion{測定される量には、実験が観測できるよりも小さいスケールが含まれるのか、もう一度説明していただけますか。\lecturetimestamp{01:05:13}}
\lecturerresponse{含まれる。それらの個々の揺らぎは分解されないが、それらの和の効果は存在する。より短いスケールを含むすべての図は、低エネルギー実験が例えば物理的な質量の二乗と呼ぶ単一の有効係数に集約される。}
\end{classroomqa}

\begin{classroomqa}
\audiencequestion{では入力項 \(m_0^2\) が他の寄与を相殺しなければならないのか。\lecturetimestamp{01:06:00}}
\lecturerresponse{測定されたスカラー質量が個々の紫外寄与よりもはるかに小さいならば、そのとおりである。要求されるこの相殺が、これから扱う微調整問題である。}
\end{classroomqa}

\section{ヒッグス階層問題}

カットオフがプランク長であると想像しよう。すると有効場の理論は、実験室のスケールから量子重力が介入しなければならないスケールまでの揺らぎを要約するよう求められることになる。素粒子物理学の単位では、
\begin{equation}
\delta\sim M_{\mathrm{Pl}}^{-1}\sim10^{-19}\,\mathrm{GeV}^{-1},
\qquad
\frac{1}{\delta^2}\sim10^{38}\,\mathrm{GeV}^2.
\label{eq:planck-cutoff}
\end{equation}
換算プランクスケールと非換算プランクスケールのどちらを正確に用いるかは、この途方もない比較には無関係である。

距離の階層も同じように際立っている。プランク長は陽子半径よりおよそ19桁小さく、議論されている加速器が直接探る距離よりおよそ16桁か17桁小さい。その隔たりを二乗することが、長さの大きな階層をスカラー質量二乗パラメータにおけるおよそ34桁の相殺に変えるのである。

ヒッグス場はスカラーであり、望ましい弱いスケールは数百GeV程度であって、\(10^{19}\) GeVではない。それにもかかわらず、式~\eqref{eq:scalar-mass-series}の項は自然にプランクスケールの大きさになり得る。それらの符号は追跡されていない。正のものも負のものもあり得て、裸の質量の二乗もどちらの符号も取り得る。しかし、電弱スケールの残りを残すには、それらの巨大な部分がおよそ34桁まで相殺しなければならない。サスキンドは最初17桁と言い、次に二乗質量の比較であるとして訂正する。この版のチューニングを測るのは二乗比である。

\begin{classroomqa}
\audiencequestion{\(\lambda\)は1より小さいか大きいか。\lecturetimestamp{01:09:50}}
\lecturerresponse{1より適度に小さいと仮定されている。逐次の冪が増大すれば摂動論は扱いにくくなり、純粋なスカラー理論も任意に大きな四重項結合で無矛盾に外挿することはできない。しかし\(\lambda\)は\(10^{-38}\)ほどには遥かに小さくない。講義では目安としてパーセントスケールの数を取る。}
\end{classroomqa}

\begin{classroomqa}
\audiencequestion{たとえ\(\lambda<1\)でも、各項が大きなカットオフ因子を運ぶため、多くの項が重要であり続けるのではないか。\lecturetimestamp{01:10:43}}
\lecturerresponse{その通りである。高次の冪は最終的に級数を抑えるが、それが起こる前に多数の項が寄与し得る。それらと裸の項がすべて共謀して、極めて小さな物理的答えを残さなければならない。}
\end{classroomqa}

ヒッグス質量パラメータが弱いスケールであると知られている理由を思い出そう。講義の規約では、対称性の破れのポテンシャルを
\begin{equation}
V(\phi)=-\frac{\mu^2}{2}\phi^2+\frac{\lambda}{4}\phi^4.
\label{eq:higgs-potential}
\end{equation}
と書く。負の二次係数は原点を不安定にし、正の四重項項はポテンシャルを安定化する。微分すると、
\begin{equation}
\frac{dV}{d\phi}=-\mu^2\phi+\lambda\phi^3
=\phi(-\mu^2+\lambda\phi^2),
\end{equation}
となり、非零の最小値
\begin{equation}
\phi_{\min}^2=f^2=\frac{\mu^2}{\lambda}.
\label{eq:higgs-minimum}
\end{equation}
を与える。黒板では四重項を\(\lambda\phi^4\)と略記するが、式~\eqref{eq:higgs-minimum}を引用する。従来の因子\(1/4\)を挿入すると、議論を変えずにそれら二行が整合する。

\(-\mu^2\)の符号は重要だが、自然さを救わない。自発的対称性の破れを駆動するのは原点における負の曲率である。繰り込みはその曲率を大きな正と負の部分でずらす。観測される破れた真空は、それらの和が単に小さいことではなく、最小値を弱いスケールに置く特定の小さな負の値を持つことを要求する。

\(W\)ボソンと\(Z\)ボソンの質量は、講義で用いられた粗い規格化において\(f\)が\(200\,\mathrm{GeV}\)程度であることを決める。\(\lambda\)が巨大でもなく途方もなく小さくもなければ、\(\mu\)は同じ一般的な程度でなければならない。したがって観測される弱いスケールの係数は、それぞれが数十桁も大きくなり得る寄与の残差である。

これはヒッグス微調整問題、階層問題、またはゲージ階層問題と呼ばれる。階層とは、通常の弱いスケールと可能な大統一スケールまたはプランクスケールとの間の途方もない比である。小さな数それ自体は逆説的である必要はない。ここで厄介なのは、極めて高い精度で相殺しなければならない、大きく一見無関係な部分から組み立てられた小さな数である。この問題は1970年代後半と1980年代前半に認識され、その相殺の背後にある機構の探求が超対称性と幾つかの競合するアイデアを動機づける。

無限の摂動級数は謎を除去するどころか鋭くする。パーセントスケールの\(\lambda\)は十分に高い冪を小さくするが、その抑制が勝つ前に多くの項が生じ得る。一ループが裸の質量を偶然に相殺するだけでは十分でない。残りの短距離過程を含めた後でも、有効な弱いスケールは安定に留まらなければならない。項の間の構造的関係は説明であろう。無関係な桁ごとの相殺は単なるチューニングにすぎない。

\section{なぜフェルミオンの質量はより良く振る舞うのか}

他の標準模型粒子は同じ種類の独立なチューニングを必要としない。ディラックフェルミオンでは、質量項は
\begin{equation}
\mathcal L_m=-m\bar\psi\psi
=-m\left(\bar\psi_L\psi_R+\bar\psi_R\psi_L\right).
\label{eq:fermion-mass}
\end{equation}
である。それは質量を含み、その二乗ではなく、左と右のカイラリティを結合する。したがって局所的な質量挿入は\(L\)を\(R\)に、または\(R\)を\(L\)に変える。

\begin{figure}[htbp]
\centering
\includegraphics[width=0.76\textwidth]{lecture_01_figure_04.png}
\caption{フェルミオン質量項と、黒板のカイラリティを変える線のスケッチ。チョークの一部が隠れているため、図は以下で代数的に再構成する。}
\lectureframe{01:23:36}
\end{figure}

講義ではしばしば handedness または helicity と言う。正確な保護はカイラリティの言葉で最も明瞭に述べられる。超相対論的粒子では helicity の描像が意図された直観を与えるが、式~\eqref{eq:fermion-mass}の左巻き場と右巻き場が運ぶ量はカイラリティである。

ゲージ相互作用はカイラリティを保存する：
\begin{equation}
L\xrightarrow{\ \gamma\ }L,
\qquad
R\xrightarrow{\ \gamma\ }R.
\label{eq:gauge-preserves-chirality}
\end{equation}
スカラーとの湯川相互作用は反対のカイラリティを結びつける：
\begin{equation}
L\xleftrightarrow{\ \phi\ }R.
\label{eq:scalar-flips-chirality}
\end{equation}
まず電子が\(m_0=0\)で始まると仮定しよう。任意個の光子を放出し再吸収しても質量を生じることはできない。なぜなら、すべてのゲージ頂点はカイラリティを保存するからである。左で始まる系列は左のままである。したがって、純粋な量子電磁力学は無質量の電子を矛盾なく許容する：その電磁的自己エネルギーはゼロから加法的な質量を生成できない。

線の描像では、左巻きカイラリティの電子から始める。最初の光子頂点でも依然として左巻きカイラリティであり、光子の伝播と再吸収の後でも依然として左巻きカイラリティである。この構成を繰り返しても、カイラリティを保存する頂点がさらに挿入されるだけである。外部2点関数は繰り込むことができるが、質量を定義する左右の構造を獲得することはできない。

スカラーの放出でそれが可能だろうか？1つのスカラー頂点は\(L\)を\(R\)に変えるが、放出されたスカラーは再吸収されなければならない。したがって、2点自己エネルギーグラフには偶数個のスカラー頂点が含まれる。2回目の反転は1回目を逆転させるので、完全なグラフは再び外部カイラリティを保存する。内部グラフがどんなに複雑になっても、放出と再吸収が頂点を対にし、始まりの質量がゼロのときに加法的な\(L\)から\(R\)への項が現れるのを妨げる。

これは反転の回数に関するパリティの議論である。1つのスカラー頂点はカイラリティを1回反転させる。すべての内部スカラー線を閉じるには第2の端点、したがって第2の反転が必要である。4つのスカラー頂点は4回反転し、以下同様である。偶数回はフェルミオンを元のカイラリティに戻す。この結論は黒板に描かれた最も単純なループに限定されない。

非ゼロの質量補正が可能になるのは、図がすでに質量挿入を含んだ後だけである。例えば、光子の放出と再吸収が1つの\(m_0\bar\psi\psi\)挿入を囲むことができる。ゲージ頂点は電気 charge の2つの冪に寄与するが、唯一のカイラリティ反転は元の質量に寄与する：
\begin{equation}
\begin{aligned}
m_e
&=m_0+c_1e^2m_0+c_2e^4m_0+\cdots\\
&=m_0\left(1+c_1e^2+c_2e^4+\cdots\right).
\end{aligned}
\label{eq:multiplicative-fermion-mass}
\end{equation}
\(m_0=0\)ならば、結果はゼロのままである。\(m_0\)が小さければ、すべての補正はその小さな数に比例する。電子が基本スケールに対してなぜ軽いのかと問うことは依然としてできるが、それを得るために莫大で無関係な数を差し引く必要はもはやない。この意味で、小さなフェルミオン質量は技術的に自然である。

質量挿入は許される系列を明示する。\(L\)から始める。光子放出は\(L\)を保つ；挿入\(m_0\)は\(L\)を\(R\)に変える；光子再吸収は\(R\)を保つ。この図は今や質量項の外部カイラリティを持つが、その値は必然的に因子\(m_0\)を含む。これが式~\eqref{eq:one-loop-scalar-mass}との決定的な対照であり、そのスカラー補正は出発点のスカラー質量をゼロに設定してもなお非零に留まる。

ゲージボソンにも類似の保護がある：ゲージ対称性は独立な質量項を禁じるので、無質量で始まるゲージボソンはすべての摂動次数で無質量のままである。標準模型では、フェルミオンと弱ゲージボソンの質量はヒッグス期待値に応答して生じる。ヒッグススケールを制御すれば、残りの軽い質量は同じ破局的調整の別個のコピーを必要としない。

これは超対称性が扱うべき問題を特定する：理論はヒッグススカラーが莫大な加法的質量補正を受けるのを防げるだろうか？答えは、適切に超対称な設定ではイエスである。なぜなら、ボソン的寄与とフェルミオン的寄与は独立ではなく関係しているからである。詳細な機構は次の講義から始まる。

\section{もう一つの微調整：真空エネルギー}

さらに深刻な繰り込み問題が1つあるが、それは重力を含めるときにのみ観測可能になる。非重力物理学では、すべてのエネルギーに定数を加えても実験は変わらない。例えば電子質量は絶対的なエネルギーではなく、真空に電子を加えるのに必要な余分なエネルギーである：
\begin{equation}
m_e=E(\text{vacuum}+\text{one electron})-E(\text{vacuum})
\end{equation}
電子の静止系において。真空とすべての励起状態が一緒にずれるならば、それらの差は変わらないままである。

それにもかかわらず、真空から真空への図は真空エネルギーをずらすのである。\(\lambda\phi^4\)理論では、一つの四次頂点の四本の脚を二つの短いループに結合する。各々の一致する伝播関数は\(1/\delta^2\)を与えるので、
\begin{equation}
\delta\rho_{\mathrm{vac}}
\sim
c_{\mathrm{vac}}\frac{\lambda}{\delta^4}.
\label{eq:vacuum-energy-divergence}
\end{equation}
より複雑な真空図は異なる結合因子と係数を持つが、次元解析は典型的にはエネルギー密度に対して同じ四次カットオフ感度を残す。

外部線の不在こそが、この補正を特定するものである。二点グラフは一粒子状態と真空を比較し、質量を変える。四点グラフは散乱を変える。入射粒子も出射粒子も持たないグラフは、真空から真空への振幅そのものを変える。一つの四次頂点が\(\lambda\)を与え、その二つの短いループがそれぞれ\(\delta^{-2}\)を与える。すると次元解析により、\(\delta^{-4}\)のエネルギー密度スケールが要求される。

重力がなければ、エネルギーの加法定数の零点を自由に選べるので、誰が気にするのかと問うこともできる。一般相対性理論はその自由を奪う。応力エネルギーが時空の曲率の源となるため、一様な真空エネルギーは空の空間を重力源にする。それは宇宙定数として現れ、観測的には暗黒エネルギーと結びついている。すると、測定された微小な真空エネルギーは、\(\delta^{-4}\)の形の寄与を生き延びなければならない。

重力を入れると問いが変わるのはこのためである。通常の量子力学では、すべての準位を同じ定数だけずらしても、すべての遷移振動数は変わらない。重力では、その定数は不活性ではない。真空の応力テンソルは計量に比例し、時空を曲げる。任意であったエネルギーの零点が、アインシュタイン方程式の源になるのである。

したがって、少なくとも2つの極端な自然さのパズルがある。講義では、真空エネルギーの調整についておよそ123桁、ヒッグス質量の調整についておよそ30--35桁という数字が挙げられている。正確な桁数は紫外スケールと規約に依存するが、この対比は疑いない。超対称性はここで議論したヒッグス問題への潜在的な応答である。超対称性が破れると、それだけでは観測された宇宙定数を説明しない。\editorialnote{これもまた、厳密に超対称なボソン--フェルミオン相殺が、観測された真空エネルギーに対する現象論的解決にならない理由である。我々の世界では超対称性は破れていなければならない。}

\section{教室での最後の質問}

\begin{classroomqa}
\audiencequestion{ヒッグス補正はカットオフ\(\delta\)に依存するが、そのカットオフは恣意的に見える。それはプランク長なのか、プランク長の十倍なのか、それとも何か別のものなのか。そのことが微調整を恣意的にするのだろうか。\lecturetimestamp{01:33:50}}
\lecturerresponse{数値的な寄与は、有効理論が信頼される最短スケールに依存する。その依存性は物理的な帳簿付けであり、逃げ道ではない。標準模型がなお意味を持つ任意のスケールに至るまでのモードは、その効果を必ず含めなければならないモードである。新しい物理が期待される距離において\(\delta\)を選ぶ---おそらくプランクスケール、あるいは\(10^{16}\,\mathrm{GeV}\)付近の大統一スケールであろう。そのスケールまでに蓄積される寄与は、すでにこの問題の下限である。同じ理論をさらに短い距離まで拡張することは、二次の感受性をより大きくするだけである。実験室とその紫外境界との間のあらゆるスケールからの寄与を何が相殺するのかを、我々は説明する義務がある。}
\end{classroomqa}

\begin{classroomqa}
\audiencequestion{スカラー粒子を定義するものは何か、またその例は何か。\lecturetimestamp{01:37:26}}
\lecturerresponse{ヒッグスがその意図された例である。スカラーは空間回転の下で異なる成分に変換しない。等価に言えば、その粒子量子はスピン0を持つ。光子はベクトル場で記述され、重力場はテンソルで記述される。この2010年の講義の時点では、素粒子のスピン0粒子はまだ検出されていなかった。\editorialnote{ヒッグスボソンは2012年にCERNで発見され、その測定された性質はスピン0と整合している。}}
\end{classroomqa}

\begin{classroomqa}
\audiencequestion{\(Z\)粒子もゲージボソンなのか。\lecturetimestamp{01:38:31}}
\lecturerresponse{はい。\(Z\)、\(W^+\)、\(W^-\)はスピン1のゲージボソンであり、光子も同様である。それらの場の方程式と相互作用は、マクスウェル的なゲージ記述の非可換あるいは電弱的一般化である。}
\end{classroomqa}

\begin{classroomqa}
\audiencequestion{なぜボソンの基底状態の寄与は、真空エネルギーにおいてフェルミオンの寄与を相殺しないのか。\lecturetimestamp{01:39:10}}
\lecturerresponse{フェルミオンループと対応するボソンループは、確かに逆の符号で入りうる。しかし、相殺は自動ではない。粒子の質量、多重度、結合が一致しなければならない。わずかな不一致でも、カットオフに敏感な大きな残差が残り、無限のダイアグラムの族が次数ごとに一致しなければならない。厳密な超対称性は、まさにそのようなボソン--フェルミオン関係を強制する。一般的な非超対称な粒子内容はそうではなく、破れた超対称性は真空エネルギーの問題を残す。}
\end{classroomqa}

\section{展望}

講義の道筋は意図的なものである。まず分子の計算により、速い変数を消去することは無視でも魔法でもないことが示される。それらの基底状態エネルギーが遅い変数に対するポテンシャルとなるのである。量子場理論はその操作を波長にわたって繰り返す。次に次元解析により、どの有効パラメータがカットオフに敏感であるかが明らかになる。スカラー質量は\(\delta^{-2}\)に比例する加法的項を受け取る。したがって小さなヒッグススケールは不自然に繊細に見える。フェルミオン質量は乗法的に繰り込まれる。なぜならカイラリティがゼロ質量極限を保護し、ゲージ対称性も同様に質量ゼロのゲージボソンを保護するからである。真空エネルギーはさらに敏感であり、\(\delta^{-4}\)としてスケールし、重力がその絶対値を観測可能にする。

超対称性はまさにこの点で講義に登場する。装飾的な新しい数学としてではなく、スカラー部門を保護された部門により近い振る舞いにする試みとしてである。次の課題は、どのような対称性がボソンとフェルミオンを結び付け、それらの相殺を組織するのに十分強く関係づけられるかを理解することである。